% Material suplementario — Artículo 2 (español). Se compila antes que manuscript.tex.
\documentclass[11pt]{article}
\newcommand{\DocLang}{es}
\input{papers/shared/preamble.tex}
\input{papers/shared/authors.tex}
\usepackage{pdflscape}
\hypersetup{pdftitle={Material suplementario: Atributos de liberacion de los lotes y reactogenicidad de VA-MENGOC-BC}}

\begin{document}
\SupplementaryNumbering
\SyntheticNotice
\begin{center}
{\Large\bfseries Material suplementario\par}
\medskip
{\large Atributos de liberación de los lotes y reactogenicidad de una vacuna de vesículas de
membrana externa: vinculación de los datos de control de calidad de \PTNLots{} lotes de VA-MENGOC-BC
con las notificaciones nacionales de eventos adversos\par}
\medskip
{\small \AuthorOne{} y \AuthorTwo\par}
\end{center}

\tableofcontents

\section{Gráficos de control}\label{sec:S-charts}
\begin{table}[htbp]
\centering\small
\caption{Señales en el gráfico de valores individuales (reglas de Nelson) y en el gráfico de media
móvil ponderada exponencialmente ($\lambda=\PTEWMALambda$) de cada atributo, en orden de producción.
Las señales indican desviaciones de un proceso estable, no no conformidad.}
\label{tab:S-charts}
\PTTableControlCharts
\end{table}

\section{Lotes nacionales y de exportación}\label{sec:S-equivalence}
\begin{table}[htbp]
\centering\small
\caption{Equivalencia de los lotes nacionales y de exportación mediante dos pruebas unilaterales en la
escala de la ventana de especificación (margen de $\pm$\PTEqMarginPct{} de la ventana): diferencia de
medias con intervalo de confianza del 90\,\% y desplazamiento de Hodges--Lehmann con intervalo de
confianza del 95\,\%.}
\label{tab:S-equivalence}
\PTTableEquivalence
\end{table}

\begin{figure}[htbp]
\centering
\includegraphics{p2_equivalence}
\caption{Diferencia entre lotes nacionales y de exportación (intervalo de confianza del 90\,\%)
respecto al margen de equivalencia (sombreado).}
\label{fig:S-equivalence}
\end{figure}

\section{Vinculación}\label{sec:S-linkage}
\begin{table}[htbp]
\centering\small
\caption{Notificaciones con VA-MENGOC-BC vinculadas con un lote liberado, por año de vacunación.}
\label{tab:S-linkage}
\PTTableLinkage
\end{table}

\begin{table}[htbp]
\centering\small
\caption{Características de las notificaciones vinculadas y no vinculadas con un lote liberado
(evaluación del sesgo de vinculación).}
\label{tab:S-linkbias}
\PTTableLinkBias
\end{table}

\section{Análisis de sensibilidad}\label{sec:S-sensitivity}
\begin{landscape}
\begin{table}[htbp]
\centering\small
\caption{Razones de odds por DE (IC 95\,\%) de las asociaciones preespecificadas en análisis
alternativos: ecuaciones de estimación generalizadas principales; restricción a enlaces exactos;
exclusión de vacunas coadministradas; conservación de los enlaces temporalmente inverosímiles; modelo
logístico bayesiano con intercepto aleatorio (intervalo de credibilidad del 95\,\%); modelo
cuasibinomial a nivel de lote de la proporción de notificaciones con la reacción.}
\label{tab:S-sensitivity}
\PTTableSensitivity
\end{table}
\end{landscape}

\section{Asociaciones exploratorias}\label{sec:S-exploratory}
{\small
\PTTableGEEExploratory{Asociaciones exploratorias entre cada atributo de los lotes (incluida la
atipicidad multivariante) y las reacciones principales y secundarias: modelos logísticos con errores
estándar robustos agrupados por lote (correlación de trabajo de independencia), ajustados como en el
análisis principal. $q$, valor de $p$ ajustado por Benjamini--Hochberg dentro de la familia
exploratoria.\label{tab:S-exploratory}}}

\end{document}
